\markdownRendererDocumentBegin
\markdownRendererWarning{The `hybrid` option has been soft-deprecated.}{The `hybrid` option has been soft-deprecated.}{Consider using one of the following better options for mixing TeX and markdown: `contentBlocks`, `rawAttribute`, `texComments`, `texMathDollars`, `texMathSingleBackslash`, and `texMathDoubleBackslash`. For more information, see the user manual at <https://witiko.github.io/markdown/>.}{Consider using one of the following better options for mixing TeX and markdown: `contentBlocks`, `rawAttribute`, `texComments`, `texMathDollars`, `texMathSingleBackslash`, and `texMathDoubleBackslash`. For more information, see the user manual at <https://witiko.github.io/markdown/>.}$\clearpage$\markdownRendererInterblockSeparator
{}\markdownRendererSectionBegin
\markdownRendererSectionBegin
\markdownRendererSectionBegin
\markdownRendererHeadingThree{Strong Connectivity}\markdownRendererParagraphSeparator
{}
\markdownRendererSectionEnd \markdownRendererSectionBegin
\markdownRendererHeadingThree{Kosaraju}\markdownRendererInterblockSeparator
{}El algoritmo de \markdownRendererStrongEmphasis{Kosaraju} encuentra las \markdownRendererStrongEmphasis{componentes fuertemente conexas (SCC)} en un grafo dirigido, donde cada componente es un conjunto de nodos que se alcanzan mutuamente. El proceso tiene dos fases:\markdownRendererHardLineBreak
{}1️⃣ Realiza un DFS normal y guarda los nodos según su orden de finalización.\markdownRendererHardLineBreak
{}2️⃣ Invierte las aristas y ejecuta DFS en ese orden inverso para formar las SCC.\markdownRendererParagraphSeparator
{}Cada nodo se visita exactamente dos veces, por lo que su complejidad es $O(n+m)$.\markdownRendererInterblockSeparator
{}\markdownRendererInputFencedCode{./_markdown_main/f05b6c5b8baa17649f25325511fd32d5.verbatim}{cpp}{cpp}\markdownRendererInterblockSeparator
{}
\markdownRendererSectionEnd \markdownRendererSectionBegin
\markdownRendererHeadingThree{2 SAT}\markdownRendererInterblockSeparator
{}
\markdownRendererSectionEnd \markdownRendererSectionBegin
\markdownRendererHeadingThree{2-SAT}\markdownRendererInterblockSeparator
{}El \markdownRendererStrongEmphasis{problema 2-SAT} consiste en determinar si es posible asignar valores de verdad a un conjunto de variables booleanas $x_1, x_2, \markdownRendererEllipsis{}, x_n$ para que una fórmula de la forma\markdownRendererHardLineBreak
{}$(a_1 \lor b_1) \land (a_2 \lor b_2) \land \dots \land (a_m \lor b_m)$ sea verdadera, donde cada $a_i$ y $b_i$ son una variable o su negación.\markdownRendererParagraphSeparator
{}Cada cláusula $(a \lor b)$ genera dos \markdownRendererStrongEmphasis{implicaciones}: ¬a → b y ¬b → a. Con esto, construimos un \markdownRendererStrongEmphasis{grafo de implicaciones} y verificamos si existe una variable $x_i$ y su negación ¬$x_i$ en la misma \markdownRendererStrongEmphasis{componente fuertemente conexa (SCC)}.\markdownRendererHardLineBreak
{}Si eso ocurre, la fórmula es \markdownRendererStrongEmphasis{insatisfacible}. Si no ocurre, la fórmula \markdownRendererStrongEmphasis{tiene solución}, que se obtiene procesando las componentes en orden topológico inverso.\markdownRendererHardLineBreak
{}Su complejidad es $O(n + m)$.\markdownRendererInterblockSeparator
{}
\markdownRendererSectionEnd \markdownRendererSectionBegin
\markdownRendererHeadingThree{2-SAT}\markdownRendererInterblockSeparator
{}El \markdownRendererStrongEmphasis{problema 2-SAT} busca determinar si se pueden asignar valores de verdad a las variables booleanas $x_1, x_2, \dots, x_n$ para que\markdownRendererHardLineBreak
{}$(a_1 \lor b_1) \land (a_2 \lor b_2) \land \dots \land (a_m \lor b_m)$ sea verdadera, donde cada $a_i$ y $b_i$ es una variable o su negación.\markdownRendererParagraphSeparator
{}Por ejemplo, la fórmula\markdownRendererHardLineBreak
{}$$(x_2 \lor \lnot x_1) \land (\lnot x_1 \lor \lnot x_2) \land (x_1 \lor x_3) \land (\lnot x_2 \lor \lnot x_3) \land (x_1 \lor x_4)$$\markdownRendererHardLineBreak
{}es verdadera cuando las variables toman los valores:\markdownRendererHardLineBreak
{}$$\markdownRendererSoftLineBreak
{}\begin{cases}\markdownRendererSoftLineBreak
{}x_1 = \text{false} \\\markdownRendererSoftLineBreak
{}x_2 = \text{false} \\\markdownRendererSoftLineBreak
{}x_3 = \text{true} \\\markdownRendererSoftLineBreak
{}x_4 = \text{true}\markdownRendererSoftLineBreak
{}\end{cases}\markdownRendererSoftLineBreak
{}$$\markdownRendererParagraphSeparator
{}Cada cláusula $(a \lor b)$ genera dos \markdownRendererStrongEmphasis{implicaciones}: ¬a → b y ¬b → a. Con esto se construye un \markdownRendererStrongEmphasis{grafo de implicaciones}, donde el problema es \markdownRendererStrongEmphasis{insatisfacible} si existe una variable $x_i$ y su negación ¬$x_i$ en la misma \markdownRendererStrongEmphasis{componente fuertemente conexa (SCC)}.\markdownRendererHardLineBreak
{}Si no ocurre, la fórmula tiene solución, y los valores se obtienen procesando las componentes en orden topológico inverso. Su complejidad es $O(n+m)$.\markdownRendererParagraphSeparator
{}$\clearpage$\markdownRendererInterblockSeparator
{}\markdownRendererInputFencedCode{./_markdown_main/d7a2049200203c4ada695b8e47fdd183.verbatim}{cpp}{cpp}\markdownRendererParagraphSeparator
{}
\markdownRendererSectionEnd \markdownRendererSectionBegin
\markdownRendererHeadingThree{Eulerian Paths and Circuits}\markdownRendererInterblockSeparator
{}Un \markdownRendererStrongEmphasis{camino Euleriano} es aquel que recorre \markdownRendererStrongEmphasis{cada arista exactamente una vez}.\markdownRendererHardLineBreak
{}Si el camino empieza y termina en el mismo nodo, se llama \markdownRendererStrongEmphasis{circuito Euleriano}.\markdownRendererParagraphSeparator
{}En \markdownRendererStrongEmphasis{grafos no dirigidos}, existe un camino Euleriano si todas las aristas están en la misma componente conexa y:\markdownRendererInterblockSeparator
{}\markdownRendererUlBeginTight
\markdownRendererUlItem Todos los nodos tienen grado par, o\markdownRendererUlItemEnd 
\markdownRendererUlItem Exactamente dos nodos tienen grado impar (estos serán los extremos del camino).\markdownRendererUlItemEnd 
\markdownRendererUlEndTight \markdownRendererInterblockSeparator
{}En \markdownRendererStrongEmphasis{grafos dirigidos}, se cumple si todas las aristas pertenecen a la misma componente y:\markdownRendererInterblockSeparator
{}\markdownRendererUlBeginTight
\markdownRendererUlItem Cada nodo tiene el mismo número de entradas y salidas (circuito Euleriano), o\markdownRendererUlItemEnd 
\markdownRendererUlItem Hay un nodo con una salida más que entradas (inicio) y otro con una entrada más que salidas (final).\markdownRendererUlItemEnd 
\markdownRendererUlEndTight \markdownRendererInterblockSeparator
{}Para \markdownRendererStrongEmphasis{construir} un circuito o camino Euleriano se usa el \markdownRendererStrongEmphasis{algoritmo de Hierholzer}, que:\markdownRendererInterblockSeparator
{}\markdownRendererOlBeginTight
\markdownRendererOlItemWithNumber{1}Empieza desde un nodo con aristas disponibles.\markdownRendererOlItemEnd 
\markdownRendererOlItemWithNumber{2}Recorre aristas sin repetir hasta regresar al inicio (subcircuito).\markdownRendererOlItemEnd 
\markdownRendererOlItemWithNumber{3}Inserta subcircuitos dentro del circuito principal hasta usar todas las aristas.\markdownRendererOlItemEnd 
\markdownRendererOlEndTight \markdownRendererInterblockSeparator
{}Su complejidad es $O(n+m)$, ya que cada arista se recorre una vez.\markdownRendererParagraphSeparator
{}$\clearpage$\markdownRendererInterblockSeparator
{}\markdownRendererInputFencedCode{./_markdown_main/675687cd69db78bd9eab9fd7ed8cd317.verbatim}{cpp}{cpp}\markdownRendererInterblockSeparator
{}
\markdownRendererSectionEnd \markdownRendererSectionBegin
\markdownRendererHeadingThree{Hamiltonian Path}\markdownRendererInterblockSeparator
{}El \markdownRendererStrongEmphasis{camino Hamiltoniano} recorre \markdownRendererStrongEmphasis{cada vértice exactamente una vez}, y si empieza y termina en el mismo vértice, se denomina \markdownRendererStrongEmphasis{circuito Hamiltoniano}.\markdownRendererHardLineBreak
{}Encontrar uno es un problema \markdownRendererStrongEmphasis{NP-hard}, ya que no existe un algoritmo eficiente general. Sin embargo, usando \markdownRendererStrongEmphasis{programación dinámica con bitmask}, se puede verificar su existencia en tiempo $O(2^n n^2)$.\markdownRendererParagraphSeparator
{}$$dp\markdownRendererWarning{Undefined link reference "u"}{Undefined link reference "u"}{Look for the text `[...][u]`.}{Look for the text `[...][u]`.}[mask][u] = \bigvee_{v \in mask,, adj\markdownRendererWarning{Undefined link reference "u"}{Undefined link reference "u"}{Look for the text `[...][u]`.}{Look for the text `[...][u]`.}[v][u]} dp\markdownRendererWarning{Undefined link reference "v"}{Undefined link reference "v"}{Look for the text `[...][v]`.}{Look for the text `[...][v]`.}[mask \setminus {u}][v]$$\markdownRendererInterblockSeparator
{}\markdownRendererInputFencedCode{./_markdown_main/dd58373418d2b4cd253efc62c51df273.verbatim}{cpp}{cpp}\markdownRendererInterblockSeparator
{}$\clearpage$\markdownRendererInterblockSeparator
{}
\markdownRendererSectionEnd \markdownRendererSectionBegin
\markdownRendererHeadingThree{De Bruijn Sequences}\markdownRendererInterblockSeparator
{}Una \markdownRendererStrongEmphasis{secuencia de De Bruijn} es una cadena que contiene \markdownRendererStrongEmphasis{todas las posibles subcadenas de longitud $n$} exactamente una vez, usando un alfabeto de $k$ símbolos.\markdownRendererHardLineBreak
{}La longitud total es $k^n + n - 1$.\markdownRendererHardLineBreak
{}Por ejemplo, con $n = 3$ y $k = 2$, una secuencia válida es \markdownRendererCodeSpan{0001011100},\markdownRendererHardLineBreak
{}que incluye todas las combinaciones binarias de 3 bits: 000, 001, 010, 011, 100, 101, 110 y 111.\markdownRendererParagraphSeparator
{}Cada secuencia De Bruijn corresponde a un \markdownRendererStrongEmphasis{camino Euleriano} en un grafo donde cada nodo representa una cadena de $n − 1$ símbolos y cada arista agrega uno nuevo.\markdownRendererHardLineBreak
{}Al recorrer el camino Euleriano y concatenar los símbolos, se obtiene la secuencia completa de longitud $k^n + n − 1$.\markdownRendererInterblockSeparator
{}\markdownRendererInputFencedCode{./_markdown_main/8687d6ccffadc698831f92145475632d.verbatim}{cpp}{cpp}\markdownRendererInterblockSeparator
{}$\clearpage$\markdownRendererInterblockSeparator
{}
\markdownRendererSectionEnd \markdownRendererSectionBegin
\markdownRendererHeadingThree{Knight’s Tour}\markdownRendererInterblockSeparator
{}Un recorrido del caballo es una secuencia de movimientos en un tablero $n × n$ donde el caballo visita cada casilla exactamente una vez.\markdownRendererSoftLineBreak
{}Si termina en el punto de inicio, se llama recorrido cerrado; de lo contrario, es abierto.\markdownRendererSoftLineBreak
{}El problema equivale a encontrar un camino Hamiltoniano en el grafo formado por las casillas como nodos y los movimientos válidos del caballo como aristas.\markdownRendererParagraphSeparator
{}El recorrido puede generarse por backtracking, pero se acelera con la regla de Warnsdorf,que elige siempre el movimiento hacia la casilla con el menor número de movimientos siguientes posibles.\markdownRendererInterblockSeparator
{}\markdownRendererInputFencedCode{./_markdown_main/6926b8b98c39940bf4786b6c87f2d68f.verbatim}{cpp}{cpp}\markdownRendererInterblockSeparator
{}$\clearpage$\markdownRendererParagraphSeparator
{}
\markdownRendererSectionEnd 
\markdownRendererSectionEnd 
\markdownRendererSectionEnd \markdownRendererSectionBegin
\markdownRendererHeadingOne{Flujos}\markdownRendererInterblockSeparator
{}\markdownRendererSectionBegin
\markdownRendererHeadingTwo{Red de Flujos}\markdownRendererInterblockSeparator
{}Una red de flujo es un grafo dirigido con:\markdownRendererInterblockSeparator
{}\markdownRendererUlBeginTight
\markdownRendererUlItem V vértices y E aristas\markdownRendererUlItemEnd 
\markdownRendererUlItem Cada arista tiene una capacidad c(u,v) ≥ 0\markdownRendererUlItemEnd 
\markdownRendererUlItem Hay una fuente s y un sumidero t\markdownRendererUlItemEnd 
\markdownRendererUlEndTight \markdownRendererInterblockSeparator
{}Un flujo f(u,v) cumple:\markdownRendererInterblockSeparator
{}\markdownRendererOlBeginTight
\markdownRendererOlItemWithNumber{1}0 $\leq$ f(u,v) $\leq$ c(u,v)\markdownRendererOlItemEnd 
\markdownRendererOlItemWithNumber{2}Conservación del flujo en todo vértice excepto s y t: suma de flujos que entran = suma de flujos que salen\markdownRendererOlItemEnd 
\markdownRendererOlEndTight \markdownRendererInterblockSeparator
{}Objetivo:\markdownRendererSoftLineBreak
{}Enviar todo el flujo posible desde s hasta t.\markdownRendererInterblockSeparator
{}
\markdownRendererSectionEnd \markdownRendererSectionBegin
\markdownRendererHeadingTwo{Grafo Residual}\markdownRendererInterblockSeparator
{}El grafo residual contiene:\markdownRendererInterblockSeparator
{}\markdownRendererUlBeginTight
\markdownRendererUlItem Aristas con capacidad disponible: residual = c(u,v) - f(u,v)\markdownRendererUlItemEnd 
\markdownRendererUlItem Aristas reversas con residual = f(u,v)\markdownRendererUlItemEnd 
\markdownRendererUlEndTight \markdownRendererInterblockSeparator
{}Si existe flujo hacia adelante, puede enviarse más;\markdownRendererSoftLineBreak
{}si existe flujo hacia atrás, puede revertirse parte del flujo.\markdownRendererParagraphSeparator
{}Este grafo cambia dinámicamente durante el algoritmo, para la lectura se tienen que realizar algunas modificaciones.\markdownRendererInterblockSeparator
{}\markdownRendererInputFencedCode{./_markdown_main/8c53e7c77894530e483479243a7b228d.verbatim}{cpp}{cpp}\markdownRendererInterblockSeparator
{}$\clearpage$\markdownRendererParagraphSeparator
{}
\markdownRendererSectionEnd \markdownRendererSectionBegin
\markdownRendererHeadingTwo{Ford-Fulkerson}\markdownRendererInterblockSeparator
{}Idea:\markdownRendererSoftLineBreak
{}Mientras exista un camino aumentante de s → t en el grafo residual:\markdownRendererInterblockSeparator
{}\markdownRendererUlBeginTight
\markdownRendererUlItem Se calcula el mínimo residual C en ese camino\markdownRendererUlItemEnd 
\markdownRendererUlItem Se aumenta el flujo en C a lo largo del camino\markdownRendererUlItemEnd 
\markdownRendererUlItem Se actualizan capacidades residuales\markdownRendererUlItemEnd 
\markdownRendererUlEndTight \markdownRendererInterblockSeparator
{}Termina cuando ya no hay camino aumentante.\markdownRendererSoftLineBreak
{}Complejidad: O(E * F) donde F = flujo máximo\markdownRendererSoftLineBreak
{}(Si las capacidades son enteras, siempre termina)\markdownRendererParagraphSeparator
{}\markdownRendererInputFencedCode{./_markdown_main/0c687b54648b4a9b296c98f8a394497b.verbatim}{cpp}{cpp}\markdownRendererInterblockSeparator
{}$\clearpage$\markdownRendererInterblockSeparator
{}
\markdownRendererSectionEnd \markdownRendererSectionBegin
\markdownRendererHeadingTwo{Edmonds-Karps}\markdownRendererInterblockSeparator
{}Mejora de Ford–Fulkerson: usar BFS para encontrar caminos aumentantes\markdownRendererSoftLineBreak
{}→ Elige siempre el camino más corto en número de aristas\markdownRendererParagraphSeparator
{}Complejidad garantizada:\markdownRendererSoftLineBreak
{}O(V * E^2) incluso con capacidades irracionales\markdownRendererParagraphSeparator
{}Es el algoritmo estándar para Max-Flow en competiciones.\markdownRendererInterblockSeparator
{}\markdownRendererInputFencedCode{./_markdown_main/ed8f0ddf5d75db801f952b1d277000da.verbatim}{cpp}{cpp}\markdownRendererInterblockSeparator
{}$\clearpage$\markdownRendererInterblockSeparator
{}
\markdownRendererSectionEnd \markdownRendererSectionBegin
\markdownRendererHeadingTwo{Teorema del Corte Minimo}\markdownRendererInterblockSeparator
{}Teorema:\markdownRendererSoftLineBreak
{}El valor del flujo máximo es igual al valor del mínimo corte\markdownRendererParagraphSeparator
{}Un s-t cut es una partición de los vértices:\markdownRendererInterblockSeparator
{}\markdownRendererUlBeginTight
\markdownRendererUlItem S contiene a s\markdownRendererUlItemEnd 
\markdownRendererUlItem T contiene a t\markdownRendererSoftLineBreak
{}Se suman capacidades de aristas de S → T\markdownRendererUlItemEnd 
\markdownRendererUlEndTight \markdownRendererInterblockSeparator
{}No se puede mandar más flujo del valor del corte mínimo\markdownRendererSoftLineBreak
{}y Ford–Fulkerson demuestra que sí se puede llegar a ese valor.\markdownRendererInterblockSeparator
{}\markdownRendererInputFencedCode{./_markdown_main/70ff454a7365de161ed57041fb72516c.verbatim}{cpp}{cpp}\markdownRendererInterblockSeparator
{}\markdownRendererSectionBegin
\markdownRendererHeadingThree{Disjoint Paths}\markdownRendererInterblockSeparator
{}Los caminos disjuntos en aristas se encuentran asignando capacidad 1 a cada arista y calculando el flujo máximo. Cada unidad de flujo representa un camino independiente entre la fuente y el sumidero.\markdownRendererParagraphSeparator
{}Para los caminos disjuntos en nodos, se divide cada vértice en dos (entrada y salida) conectados por una arista con capacidad 1, asegurando que cada nodo solo pueda usarse en un camino. El máximo flujo en esta red equivale al número máximo de caminos nodo-disjuntos.\markdownRendererInterblockSeparator
{}\markdownRendererInputFencedCode{./_markdown_main/12b83366a1243cd2a5c334cb183814e0.verbatim}{cpp}{cpp}\markdownRendererParagraphSeparator
{}
\markdownRendererSectionEnd 
\markdownRendererSectionEnd \markdownRendererSectionBegin
\markdownRendererHeadingTwo{Bipartite Matching}\markdownRendererInterblockSeparator
{}El emparejamiento máximo bipartito se reduce fácilmente a un problema de flujo máximo: se conectan todos los nodos izquierdos con la fuente y todos los derechos con el sumidero, asignando capacidad 1 a cada arista. El valor del flujo máximo equivale al tamaño del emparejamiento máximo.\markdownRendererParagraphSeparator
{}El teorema de Hall establece cuándo existe un emparejamiento perfecto: para todo subconjunto de nodos izquierdos X, su conjunto de vecinos f(X) debe cumplir |X| ≤ |f(X)|. Si alguna elección viola esta condición, no puede formarse un emparejamiento perfecto.\markdownRendererInterblockSeparator
{}\markdownRendererInputFencedCode{./_markdown_main/12cc1933ec32457733a6ff2f55866399.verbatim}{cpp}{cpp}\markdownRendererInterblockSeparator
{}$\clearpage$\markdownRendererInterblockSeparator
{}
\markdownRendererSectionEnd \markdownRendererSectionBegin
\markdownRendererHeadingTwo{Puentes}\markdownRendererInterblockSeparator
{}Un puente en un grafo es una arista cuya eliminación dividiría el grafo en dos o más componentes conexas adicionales.\markdownRendererInterblockSeparator
{}\markdownRendererInputFencedCode{./_markdown_main/4bc6fdbadbeadc58057906e2acb814ee.verbatim}{cpp}{cpp}\markdownRendererInterblockSeparator
{}
\markdownRendererSectionEnd \markdownRendererSectionBegin
\markdownRendererHeadingTwo{Puntos de articulacion}\markdownRendererInterblockSeparator
{}\markdownRendererInputFencedCode{./_markdown_main/40b24ad61f658a2dcdd3d8827acafc9d.verbatim}{cpp}{cpp}
\markdownRendererSectionEnd 
\markdownRendererSectionEnd \markdownRendererDocumentEnd